\chapter*{Presentación}
\addcontentsline{toc}{chapter}{Presentación}
\markboth{Presentación}{}
El segundo volumen de esta colección estudia cómo medir conjuntos, integrar funciones y convertir la integración en una herramienta para construir espacios completos y representar funcionales. El recorrido comienza con operaciones de conjuntos y termina con los teoremas de Radon--Nikodym y de representación de Riesz. Las seis partes y los dieciséis capítulos organizan las 33 lecturas de \emph{Análisis Matemático 2}, de Miguel Ángel Valencia Bucio \cite{valencia2}, sin identificar un capítulo del libro con una sesión del curso.

Eder Said Martinez Ramirez es el autor de esta edición: su organización, ampliación explicativa, desarrollo de soluciones y correcciones. Las lecturas de Valencia Bucio constituyen la fuente del curso y conservan ese crédito; las contribuciones didácticas externas se identifican localmente y en las referencias. En particular, se incorporan explicaciones inspiradas en materiales de Egor Maximenko sobre Carathéodory, modos de convergencia, errores de cuantificadores y completitud de espacios de funciones. La atribución de un material no sustituye la verificación de sus argumentos: cuando una fuente contiene una errata o un enunciado falso, se explican el problema y la corrección.

Los tratados de Rudin, Folland, Halmos y Bogachev ofrecen un marco de consulta y profundización para la integración, los espacios de funciones y la representación de medidas \cite{rudin1987,folland1999,halmos1950,bogachev2007}. Las referencias finales identifican sus ediciones reales y distinguen las fechas originales de las reimpresiones.

\section*{Prerrequisitos y modo de lectura}
Se presupone el contenido básico del primer volumen: espacios métricos, continuidad, completitud, espacios normados, operadores y geometría de Hilbert. También se utilizan cálculo de una variable y álgebra lineal. Las operaciones con familias numerables, los números reales extendidos, los cocientes por igualdad casi en todas partes y los cambios de cuantificadores se revisan cuando aparecen. En los puntos donde interviene el axioma de elección se distingue la existencia abstracta de una construcción explícita.

Antes de aplicar un teorema de convergencia, identifique cuatro datos: qué convergencia se conoce, cuál se busca, qué controla los valores grandes de las funciones y qué impide que la masa escape hacia conjuntos cada vez más lejanos. En una prueba, compruebe primero que los conjuntos y las funciones utilizados son medibles. Finalmente, si aparece una resta de integrales, verifique que no se está formando la expresión indeterminada $\infty-\infty$.

Los ejercicios conservan su correspondencia con la lectura y el número original en el registro de cobertura que acompaña las fuentes LaTeX. Las soluciones explican los incisos y señalan los casos en que es necesario modificar una hipótesis. Las ampliaciones se identifican por su procedencia. El índice general permite recorrer la arquitectura temática; el índice alfabético y el de notación permiten consultar conceptos concretos. La correspondencia final hace visible la relación entre esta organización y las fuentes.

\section*{Créditos de diseño y alcance}
El diseño toma como referencia visual los \emph{Apuntes de Homotopía Estable} del repositorio público de MrPaborru (\url{https://github.com/MrPaborru/Apuntes-Homotop-a-Estable}), combinados con una composición tradicional de libros de matemáticas. El emblema de la ESFM procede de su sitio oficial. Su uso identifica la adscripción académica del autor; esta colección es una edición personal, no una publicación oficial ni una obra avalada por el IPN. Se conserva una paleta inspirada en los colores de la ESFM y una tipografía Latin Modern. Los detalles de procedencia gráfica se documentan en las fuentes LaTeX.

El alcance de la resolución de ejercicios son las listas propuestas en las 33 lecturas. Los dos exámenes conservados junto con ellas se consultaron como material de contraste, pero no forman una lista adicional de problemas resueltos de esta edición. Los libros de la bibliografía ofrecen contexto y profundización; no se atribuyen a sus autores páginas, enunciados o citas textuales que no se hayan identificado.

\section*{Convenciones}
Escribimos $\N=\{1,2,\ldots\}$, $\F=\R$ o $\C$ y $\ind A$ para la función indicadora de $A$. Una medida sin otra calificación es no negativa. Adoptamos $0\cdot\infty=0$ al integrar funciones no negativas, pero no definimos $\infty-\infty$. Una función integrable es medible y tiene integral del valor absoluto finita; cuando se permite una integral extendida se indica expresamente.

Una propiedad se cumple casi en todas partes cuando falla solamente en un conjunto contenido en un conjunto medible nulo; los representantes que se construyen en las pruebas se eligen medibles. En una medida no completa no se puede modificar arbitrariamente una función sobre un subconjunto no medible de un conjunto nulo y conservar automáticamente su medibilidad. Los espacios $L^p$ son espacios de clases de equivalencia, aunque por comodidad se escriba $f$ en lugar de $[f]$ después de justificar que la operación no depende del representante.
